\section{Onde aprofundar}

Se for ler um só livro, comece pelo de Kent Beck (TDD) ou pelo
\emph{Accelerate} (visão de entrega). Os livros de TDD e de testes unitários
estão na tabela do capítulo~\ref{cap:aula1}.

\subsection*{Livros}
\small
\begin{longtable}{@{}p{1.8cm} p{4.4cm} p{2.6cm} p{3.8cm}@{}}
\toprule
Tema & Livro & Autor(es), ano & Por que ler \\
\midrule \endhead
Carga & \citetitle{livro:gregg2020}, 2ª ed. & \citeauthor{livro:gregg2020},
  \citeyear{livro:gregg2020} & Métricas, saturação e análise de desempenho \\
Resiliência & \citetitle{livro:nygard2018}, 2ª ed.
  & \citeauthor{livro:nygard2018}, \citeyear{livro:nygard2018}
  & Circuit breaker, bulkhead, timeouts e falhas em cascata \\
Caos & \citetitle{livro:rosenthal2020} & \citeauthor{livro:rosenthal2020},
  \citeyear{livro:rosenthal2020} & Os princípios contados por quem criou a
  prática na Netflix \\
Entrega & \citetitle{livro:forsgren2018} & \citeauthor{livro:forsgren2018},
  \citeyear{livro:forsgren2018} & As métricas DORA: medir resultado, não
  atividade \\
\bottomrule
\end{longtable}
\normalsize
